\question[4]
Gegeben ist die folgende Python-Funktion und ihr Aufruf:

\begin{lstlisting}
def rechteck_flaeche(breite, hoehe):      # Zeile 1
    flaeche = breite * hoehe              # Zeile 2
    return flaeche                        # Zeile 3

ergebnis = rechteck_flaeche(4, 5)         # Zeile 4
\end{lstlisting}

Ordne die folgenden Fachbegriffe den passenden Stellen im Code zu:
\texttt{Name, Kopf, Rumpf, Parameter, Argument, Aufruf, Rückgabewert}

\begin{solutionbox}{5cm}
\begin{lstlisting}[language={}]
Name           : rechteck_flaeche
Kopf           : Zeile 1 (def rechteck_flaeche(breite, hoehe):)
Rumpf          : Zeile 2-3 (der eingerückte Code der Funktion)
Parameter      : breite, hoehe (in der Definition, Zeile 1)
Argument       : 4, 5 (die konkreten Werte beim Aufruf, Zeile 4)
Aufruf         : Zeile 4 (rechteck_flaeche(4, 5))
Rückgabewert   : flaeche, also 20 (der Wert nach return)
\end{lstlisting}
\end{solutionbox}